% ID: SEL-D0016-BD-E07
% ESTADO: cotejada. Segunda pasada visual completa y borrador compilado revisados.
% Fuente: Selectividad/matematicas_aplicadas/matematicas_aplicadas_P_D.pdf
% SHA-256: 11865b83e509d9cc2723c752f7a81f714378e1b651f4a4d010dab4a4576b74c2
% Página: 2; bloque D; ejercicio original 7.

La distancia en kilómetros recorrida al día por los vehículos de una
empresa de coches de alquiler sigue una distribución Normal de media
desconocida y varianza 225. Se toma una muestra aleatoria simple de 36 coches
y se obtiene el intervalo de confianza $(153.65,162.35)$ para la media poblacional.
\begin{enumerate}[label=\alph*)]
  \item \textbf{(1 punto)} Calcule la media muestral y el error máximo de
  estimación para ese intervalo de confianza.
  \item \textbf{(0.5 puntos)} Si con el mismo nivel de confianza, aumentamos
  el tamaño muestral, \textquestiondown{}cómo se vería afectado el error?
  \item \textbf{(1 punto)} Con un nivel de confianza del 95\%,
  \textquestiondown{}cuál debería ser el tamaño mínimo de la muestra para
  que el error cometido sea inferior a $3\,km$?
\end{enumerate}
